

\newpage

\paragraph{Problem 2 answer}~\GradescopeComment{This should be on page 7 (top) of your PDF.}

%%%%%%%%%%%%%%%%%%%%%%%%%%%%%%%%%%%%%%%%%%%%%%%%%%%%%% 
%%%%%%%%%%%%%%%%%%%%%%%%%%%%%%%%%%%%%%%%%%%%%%%%%%%%%% 
%% 
%% PROBLEM 2(a) ANSWER 
%% 
%%%%%%%%%%%%%%%%%%%%%%%%%%%%%%%%%%%%%%%%%%%%%%%%%%%%%% 
%%%%%%%%%%%%%%%%%%%%%%%%%%%%%%%%%%%%%%%%%%%%%%%%%%%%%% 

\noindent\begin{answerBox}{3.5in}   % please don't edit this line, doing so will throw off the page layout 

\paragraph{(a)}  

% Define your reduction $r$
% by defining, for every possible instance $(x=(x_1, x_2, \ldots, x_n), T)$ of  \prob{Subset Sum},
% what output $r(x, T)$ the reduction gives.
% The output will be a tuple $(G=(V, E), s, t, k)$.
% Your definition must state precisely how $G$, $s$, $t$, and $k$ are defined for the given $(x, T)$.
% \emph{Hint: Construct $G$ so its $s$-$t$ paths correspond to the subsets $S\subseteq\{1,\ldots,n\}$.}

Fix an arbitrary \prob{Subset Sum} instance $(x=(x_1, x_2, \ldots, x_n), T)$.
The output of the reduction $r$ on input $(x, T)$ is $r(x, T)=(G, s, t, k)$ where $G$, $s$, $t$ and $k$ are defined as follows.

\bigskip 

\begin{blue}

  \REPLACEME
  
\end{blue}

\end{answerBox}

%%%%%%%%%%%%%%%%%%%%%%%%%%%%%%%%%%%%%%%%%%%%%%%%%%%%%% 
%%%%%%%%%%%%%%%%%%%%%%%%%%%%%%%%%%%%%%%%%%%%%%%%%%%%%% 
%% 
%% PROBLEM 2(b) ANSWER 
%% 
%%%%%%%%%%%%%%%%%%%%%%%%%%%%%%%%%%%%%%%%%%%%%%%%%%%%%% 
%%%%%%%%%%%%%%%%%%%%%%%%%%%%%%%%%%%%%%%%%%%%%%%%%%%%%% 

\begin{answerBox}{2in}   % please don't edit this line, doing so will throw off the page layout 

\paragraph{(b)}  

% Show that your reduction is a \emph{polynomial-time} reduction,
% by explaining carefully how, given $(x, T)$, one can compute the output $r(x, T)=(G, s, t, k)$
% in time polynomial in the input size $\Theta(n\log_2 T)$.

\begin{blue}

  \REPLACEME
  
\end{blue}

\end{answerBox}

%%%%%%%%%%%%%%%%%%%%%%%%%%%%%%%%%%%%%%%%%%%%%%%%%%%%%% 
%%%%%%%%%%%%%%%%%%%%%%%%%%%%%%%%%%%%%%%%%%%%%%%%%%%%%% 
%% 
%% PROBLEM 2(c) ANSWER 
%% 
%%%%%%%%%%%%%%%%%%%%%%%%%%%%%%%%%%%%%%%%%%%%%%%%%%%%%% 
%%%%%%%%%%%%%%%%%%%%%%%%%%%%%%%%%%%%%%%%%%%%%%%%%%%%%% 

\begin{answerBox}{3.2in}   % please don't edit this line, doing so will throw off the page layout 

\paragraph{(c)}

% Consider the example input $(x, T)$ with $x=(5, 9, 10, 20)$ and $T=29$.
% Calculate the output $(G, s, t, k) = r(x, T)$ for this input.
% Draw the resulting edge-weighted digraph.  What is $k$? 

For the input $(x, T)$ with $x=(5, 9, 10, 20)$ and $T=29$,
the output $r(x, T)$ of the reduction is as follows.
% 
The value of $k$ is
\BLUE{
  \REPLACEME
.}

Here is a drawing of the graph $G$:
\smallskip

\begin{center}
\begin{blue}

  \REPLACEME
  
    %   \begin{tikzpicture}[
    %       vertex/.style={circle, inner sep=0pt, minimum width=15pt, draw}, 
    %       label/.style={pos=0.5, fill=white, inner sep=2pt}, 
    %       every path/.style={->}
    %       ]

    %     %% NODES 
        
    %     \node[vertex] (v0) {$v_0$}; 
    %     \node[vertex] (v1) [right =of v0, yshift=0pt] {$v_1$}; 
    %     \node[vertex] (v2) [right =of v1, yshift=0pt] {$v_2$}; 

    %     %% EDGES AND EDGE WEIGHTS

    %     \path (v0) edge[bend left=0] node [label] {1} (v1); 
    %     \path (v1) edge[bend left=0] node [label] {2} (v2); 

    %     \path (v0) edge[bend left=30] node [label] {3} (v2); 

    %   \end{tikzpicture}

    % Or you can delete the block above, and instead draw the forest by hand,
    % scan it to a PDF, then include the PDF in your LaTeX document
    % using these commands:

    %   \includegraphics[height=2.5in,width=6in]{YOUR_FILE_NAME_HERE}   % FILE NAME WITHOUT THE .PDF EXTENSION

    % See https://texblog.org/2012/02/23/crop-figures-with-includegraphics
    % if you want to crop the included image.

\end{blue}
\end{center}

\end{answerBox}

%%%%%%%%%%%%%%%%%%%%%%%%%%%%%%%%%%%%%%%%%%%%%%%%%%%%%% 
%%%%%%%%%%%%%%%%%%%%%%%%%%%%%%%%%%%%%%%%%%%%%%%%%%%%%% 
%% 
%% PROBLEM 2(d) ANSWER 
%% 
%%%%%%%%%%%%%%%%%%%%%%%%%%%%%%%%%%%%%%%%%%%%%%%%%%%%%% 
%%%%%%%%%%%%%%%%%%%%%%%%%%%%%%%%%%%%%%%%%%%%%%%%%%%%%% 

\newpage

\paragraph{Problem 2 answer, continued}~\GradescopeComment{This should be on page 8 (top) of your PDF.}

\noindent\begin{answerBox}{2.8in}   % please don't edit this line, doing so will throw off the page layout
  
\paragraph{(d)}

% Show the ``if'' direction.
% Given an arbitrary \prob{Subset Sum} instance $(x, T)$,
% let $G$, $s$, $t$, and $k$ be the output of your reduction $r$ given $(x, T)$.  That is, $(G, s, t, k) = r(x, T)$.
% Describe carefully how, given any $s$-$t$ path $P$ in $G$ with $\weight P = k$,
% one could obtain a subset $S\subseteq \{1,\ldots,n\}$ such that $\sum_{i\in S} x_i = T$.

Fix an arbitrary \prob{Subset Sum} instance $(x=(x_1, x_2, \ldots, x_n), T)$.
% [review 2026-09-27] fix: K: reduction output written r(I) but the instance is (x, T) -> r(x, T)
Let $G=(V,E)$, $s, t$, and $k$ be the output $r(x, T)$ of the reduction $r$ given $(x, T)$.
Assume that $G$ has an $s$-$t$ path of weight $k$.
Let $P$ be such a path.
From $P$, a subset $S\subseteq \{1,\ldots, n\}$ such that $\sum_{i\in S} x_i = T$ can be obtained as follows.
    
\bigskip 

\begin{blue}
  \REPLACEME
\end{blue}

\bigskip 

\end{answerBox}

%%%%%%%%%%%%%%%%%%%%%%%%%%%%%%%%%%%%%%%%%%%%%%%%%%%%%% 
%%%%%%%%%%%%%%%%%%%%%%%%%%%%%%%%%%%%%%%%%%%%%%%%%%%%%% 
%% 
%% PROBLEM 2(e) ANSWER 
%% 
%%%%%%%%%%%%%%%%%%%%%%%%%%%%%%%%%%%%%%%%%%%%%%%%%%%%%% 
%%%%%%%%%%%%%%%%%%%%%%%%%%%%%%%%%%%%%%%%%%%%%%%%%%%%%% 

\begin{answerBox}{1.5in}   % please don't edit this line, doing so will throw off the page layout 
  
\paragraph{(e)}

% For the input $(x, T)$ from Part (b), the answer is \textsf{yes},
% so, for $(G, s, t, k)=r(x, T)$, there should be an $s$-$t$ path in $G$ of weight $k$.
% Draw such a path $P$ in $G$.
% What is the subset $S\subseteq \{1,2,3,4\}$ corresponding to the path $P$?

The value of $k$ is
\BLUE{$
  \REPLACEME
  $.}

An $s$-$t$ path of weight $k$ is \BLUE{\(P=(
  \REPLACEME    % A comma-separated list of vertices.
  )\)}

The corresponding subset $S$ is \BLUE{\(S=\{
  \REPLACEME                    % A comma-separated list of indices.
\}\)}

\end{answerBox}

%%%%%%%%%%%%%%%%%%%%%%%%%%%%%%%%%%%%%%%%%%%%%%%%%%%%%% 
%%%%%%%%%%%%%%%%%%%%%%%%%%%%%%%%%%%%%%%%%%%%%%%%%%%%%% 
%% 
%% PROBLEM 2(f) ANSWER 
%% 
%%%%%%%%%%%%%%%%%%%%%%%%%%%%%%%%%%%%%%%%%%%%%%%%%%%%%% 
%%%%%%%%%%%%%%%%%%%%%%%%%%%%%%%%%%%%%%%%%%%%%%%%%%%%%% 

\begin{answerBox}{2.5in}   % please don't edit this line, doing so will throw off the page layout 

\paragraph{(f)} 

% Now show the ``only if'' direction.
% Given an arbitrary \prob{Subset Sum} instance $(x, T)$,
% let $G=(V, E)$, $s$, $t$, and $k$ be the output of your reduction $r$ given $(x, T)$.  That is, $(G, s, t, k) = r(x, T)$.
% Describe carefully how, given any subset $S\subseteq \{1,\ldots,n\}$ such that $\sum_{i\in S} x_i = T$,
% one could obtain from $S$ an $s$-$t$ path $P$ in $G$ with $\weight P = k$.

Fix an arbitrary \prob{Subset Sum} instance $(x=(x_1,\ldots, x_n), T)$.
Let $G=(V,E)$, $s, t$, and $k$ be the output $r(x, T)$ of the reduction $r$ given $(x, T)$.
Assume that there is a subset $S\subseteq \{1,\ldots,n\}$ such that $\sum_{i\in S} x_i = T$.
From $S$, an $s$-$t$ path $P$ with $\weight P = k$ can be obtained as follows.

\bigskip 

\begin{blue}
  \REPLACEME 
\end{blue}

\end{answerBox}

%%%%%%%%%%%%%%%%%%%%%%%%%%%%%%%%%%%%%%%%%%%%%%%%%%%%%% 
%%%%%%%%%%%%%%%%%%%%%%%%%%%%%%%%%%%%%%%%%%%%%%%%%%%%%% 
%% 
%% PROBLEM 2 COLLABORATORS AND CITATIONS 
%% 
%%%%%%%%%%%%%%%%%%%%%%%%%%%%%%%%%%%%%%%%%%%%%%%%%%%%%% 
%%%%%%%%%%%%%%%%%%%%%%%%%%%%%%%%%%%%%%%%%%%%%%%%%%%%%% 

\vfill

\begin{answerBox}{2in}          % please don't edit this line, doing so will throw off the page layout
  
  \paragraph{Problem \arabic{problem} collaboration and citations}~\GradescopeComment
  {This should be on page 8 (bottom) of your PDF.}

  I collaborated on this problem with the following people:
  \begin{blue}
    \begin{itemize}
    \item \REPLACEME                    % write "none" if none
    \end{itemize}
  \end{blue}

  My answers use ideas from the following sources (other than the lecture notes and textbooks): 
  \begin{blue}
    \begin{itemize}
    \item \REPLACEME                    % write "none" if none
    \end{itemize}
  \end{blue}

\end{answerBox}


%%% Local Variables:
%%% mode: latex
%%% TeX-master: "homework_11"
%%% End:
